\chapter*{Введение}
\addcontentsline{toc}{chapter}{Введение}

% =========================================================
% Рабочая структура введения.
% Это НЕ финальная редакция, а удобный каркас для переноса и правки.
% Во всех блоках сохранены указания, откуда брать исходный материал:
%   - ВКР
%   - презентация
% При последующей шлифовке служебные комментарии можно удалять.
% =========================================================

% Общая опора:
% ВКР: введение, с. 4--11
% Презентация: слайды 2--10, 39--43

\section*{Актуальность темы исследования}
% Опора:
% ВКР: введение, с. 4--6
% Презентация: слайды 3--5
%
% Здесь показать:
% - трудоёмкость баллистического проектирования межпланетных перелётов с ЭРДУ;
% - необходимость анализа больших массивов траекторий;
% - отсутствие универсальных аналитических решений;
% - практическую потребность в компактных формулах для ускорения проектирования.
%
% Черновой текст:

\section*{Степень научной разработанности темы исследования}
% Опора:
% ВКР: введение, с. 5--9
% Презентация: слайд 6
%
% Здесь кратко собрать:
% - аналитические подходы;
% - численные методы оптимизации;
% - методы продолжения / исследования ветвления решений;
% - существующие подходы к приближённым формулам;
% - место символьной регрессии в этом ряду.
%
% Черновой текст:

\section*{Проблема и мотивация предлагаемого подхода}
% Опора:
% ВКР: введение, с. 8--11
% Презентация: слайды 7--8
%
% Здесь зафиксировать:
% - почему трудно аппроксимировать реальные массивы траекторий;
% - почему вводится модельное многообразие;
% - почему сначала изучается упрощённая постановка, а затем результаты
%   используются в более прикладной задаче.
%
% Черновой текст:

\section*{Цель и задачи исследования}
% Опора:
% ВКР: введение, с. 8--9
% Презентация: слайд 9
%
% Цель:
%
% Задачи:
% 1.
% 2.
% 3.
% 4.

\section*{Объект и предмет исследования}
% Опора:
% ВКР: введение, с. 4--11
% Презентация: слайды 7--9
%
% Объект исследования:
%
% Предмет исследования:

\section*{Научная новизна}
% Опора:
% ВКР: введение, с. 8--9
% Презентация: слайд 10
%
% Удобно оформить нумерованным списком:
% 1.
% 2.
% 3.
% 4.
%
% Черновой текст:

\section*{Теоретическая значимость исследования}
% Опора:
% ВКР: введение, с. 9; заключение, с. 71
% Презентация: слайды 20--25, 27--32
%
% Здесь подчеркнуть:
% - исследование структуры модельного многообразия;
% - выделение точек оптимальной дальности;
% - построение аппроксимированного многообразия;
% - связь плоской и пространственной постановок через $R$-преобразование.
%
% Черновой текст:

\section*{Практическая значимость исследования}
% Опора:
% ВКР: введение, с. 9; глава 5, с. 60--70
% Презентация: слайды 34--38
%
% Здесь обязательно учесть:
% - алгоритм выбора начального приближения;
% - сценарии Земля--Марс и Земля--Венера;
% - программную реализацию ILTT\_Toolbox;
% - зарегистрированную программу (РИД) ILTT\_Toolbox.
%
% Черновой текст:

\section*{Методология и методы исследования}
% Опора:
% ВКР: главы 1--5, с. 12--70
% Презентация: слайды 12--17, 34--38
%
% Здесь собрать:
% - математическое моделирование движения КА;
% - принцип максимума Понтрягина;
% - решение двухточечных краевых задач;
% - дискретное продолжение;
% - символьную регрессию;
% - вычислительные эксперименты;
% - оценку качества начального приближения.
%
% Черновой текст:

\section*{Степень достоверности результатов}
% Опора:
% ВКР: главы 2--5, с. 15--70
% Презентация: слайды 17, 24--25, 35--38
%
% Здесь удобно отразить:
% - строгую математическую постановку;
% - согласованность с физическими ограничениями;
% - численное построение многообразий;
% - проверку аппроксимаций на области параметров и в предельных случаях;
% - массовые вычислительные эксперименты.
%
% Черновой текст:

\section*{Положения, выносимые на защиту}
% Опора:
% ВКР: введение, с. 8--9
% Презентация: слайд 39
%
% Удобно оформить кратким нумерованным списком:
% 1.
% 2.
% 3.
% 4.

\section*{Апробация результатов}
% Опора:
% ВКР: введение, с. 9--11
% Презентация: слайды 42--43
%
% Здесь перечислить конференции и семинары в согласованной форме.
%
% Черновой текст:

\section*{Публикации по теме диссертации}
% Опора:
% ВКР: публикации автора и список литературы, с. 76--79
% Презентация: слайды 40--41
%
% Здесь потом развести:
% - статьи и препринты по теме;
% - материалы конференций;
% - при необходимости отдельно отметить публикации, входящие в перечень ВАК.
%
% Черновой текст:

\section*{Личный вклад автора}
% Опора:
% ВКР: введение, с. 9
% Презентация: общая структура доклада, слайды 1--39
%
% Здесь зафиксировать:
% - что именно выполнено лично автором;
% - что сделано в совместных работах;
% - где роль автора должна быть явно оговорена.
%
% Черновой текст:

\section*{Соответствие диссертации паспорту специальности 1.1.7}
% Опора:
% ВКР: титульные данные, введение, с. 1, 8--9
% Презентация: слайды 1, 9--10
%
% Здесь явно показать:
% - какие разделы теоретической механики затрагиваются;
% - почему работа относится к специальности 1.1.7;
% - почему работа относится к физико-математической отрасли наук.
%
% Черновой текст:

\section*{Объём и структура диссертации}
% Опора:
% ВКР: оглавление, с. 2
% Презентация: слайд 2
%
% Ниже оставлен служебный блок шаблона, который автоматически подтянет
% количество глав, приложений, рисунков, таблиц и источников.

\textbf{Объём и структура диссертации.}
Диссертация состоит из~введения,
\formbytotal{totalchapter}{глав}{ы}{}{},
заключения и
\formbytotal{totalappendix}{приложен}{ия}{ий}{}.
Полный объём диссертации составляет
\formbytotal{TotPages}{страниц}{у}{ы}{}, включая
\formbytotal{totalcount@figure}{рисун}{ок}{ка}{ков} и
\formbytotal{totalcount@table}{таблиц}{у}{ы}{}.
Список литературы содержит
\formbytotal{citenum}{наименован}{ие}{ия}{ий}.
